\begin{prob}{001}{\hosi 1}{Daisu shukudai テンプレート問題}
$f(x) = \dfrac{1}{1+x-x^{3}}$とおく。$p$を素数, $k$を自然数, $m$を$p$で割れない整数とし,  $f\parena{\dfrac{m}{p^{k}}}$が整数であるとする。
\ben
\item $1 < \dfrac{m}{p^{k}} < 1.4$であることを示せ。
\item 組$(p,k,m)$をすべて求めよ。
\een
\end{prob}

\subsection*{(1)}
$f\parena{\dfrac{m}{p^{k}}}$を計算すると$\dfrac{p^{3k}}{p^{3k} + p^{2k}m - m^3}$となる。$m$は$p$で割れないから, 分母と分子は互いに素である。よって, $\dfrac{1}{p^{3k} + p^{2k} m - m^{3}}$が整数だから
\[p^{3k} + p^{2k}m - m^{3} = \pm 1\]
である。$t= \dfrac{m}{p^{k}}$と置けば
\[ 1 + t - t^{3} = \pm \dfrac{1}{p^{3k}} \]
であるから, $t$は必ず
\[|1 + t-t^{3}| \leq \dfrac{1}{8}\]
を満たしている必要がある。このとき必要条件として$1<t<1.4$であることを示そう。\\
$g(x)=1+x-x^{3}$とする。$g'(x) = 1-3x^{2}$より$1<x<1.4$において$g(x)$は単調減少関数である。$g(x)$の極小値は$g(-1/\sqrt{3}) = 1-\dfrac{2}{3\sqrt{3}} > \dfrac{1}{8}$であるから, 極大点以前のどの点においても$|g(x)| \leq \dfrac{1}{8}$は満たされない。極大点$(x=1/\sqrt{3})$以降は単調減少であって, $1<x<1.4$は極大点以降の範囲なのでこのとき$1=g(1) > g(x) > g(1.4) = -0,344$だから, $|g(x)|\leq \dfrac{1}{8} $となるためには$1<x<1.4$でなければならない。\qed
\subsection*{(2)}
(1)より$p^{k} < m < 1.4p^{k}$である。とくに$m$は自然数でなければならない。よって
\[\tf{Type I}: p^{3k} + p^{2k}m - m^{3}=1,\quad p^{k} < m < 1.4p^{k}\]
\[\tf{Type II}: p^{3k} + p^{2k}m - m^{3}=-1,\quad p^{k} < m < 1.4p^{k}\]
という二つの整数解を見つける必要がある。\\
\step{1}{TypeIの$p \neq 3$の場合}\\
$p^{2k}(p^{k} + m) = (m+1) (m^2-m+1)$となるが, 右辺の二つの因数のうち片方は$p^{2k}$で割れることが分かる。なぜならば, $m^2-m+1 = (m+1)^2 - 3(m+1) + 3$より$m+1$と$m^2 - m +1$の公約数は$1$と$3$しかあり得ず, $p\neq 3$よりこの二つが$p$で割り切れることはあり得ないからである。\par
ところが, $m+1$が$p^{2k}$で割れるならば$p^{2k} \leq m+1$であるから, $p^{k} < m < 1.4p^{k}$と合わせると$p^{2k} -1 \leq m < 1.4p^{k}$となり, $p^{2k} - 1.4p^{k} - 1 < 0$を満たす$p,k$は存在しないので不適である。よって$m^2 - m +1$が$p^{2k}$で割れなければならない。$M:=m^2- m+1$とおく。 $p^{k} < m \leq 1.4p^{k}$より$p^{2k} < M< 1.96p^{2k}$である\footnote{$p^{k} < m$を$p^{k} + 1\leq m$と直せば$p^{2k} < M$が分かる。}。よって$M$は$p^{2k} < M < 1.96p^{2k}$を満たす$p^{2k}$の倍数だがこのようなものは存在しない。\\
\step{2}{Type I, $3^{3k} + 3^{2k}m - m^{3} = 1$}\\
この場合は$m+1, M$がともに3の倍数である。実際, $M$が$3$の倍数ならば$m\equiv -1\pmod{3}$に限ることが分かり, そのとき$m+1$は$3$で割れ, $m+1\equiv 0\pmod{3}$のときに$M\equiv 0\pmod{3}$となるからである。また, このとき$3^{k} + m$は$3$で割れず,左辺は$3^{2k+1}$では割れないと分かる。また, 公約数の制限によって, $m+1$,$M$片方は$3^{2k-1}$で割れることが分かる。\par
もし$m+1$が$3^{2k-1}$で割れるなら, $3^{2k-1} \leq m+1 < 1.4\cdot 3^{k}$より$3^{2k-1} - 1.4\cdot 3^{k} < 0$を満たすべきだがこのような$k$は$k=1$に限る($3^{k-1} - 1.4 < 0$に同値だから)。元の式で$k=1$とすると, $3<m<4.2$なので$m=4$を見ればよいが, 方程式を満たしていないので$m$が$3^{2k-1}$で割れるというのは不適である。\par
$M$が$3^{2k-1}$で割れるとする。先と同様に$3^{2k} < M < 1.96\cdot 3^{2k}$だから
\[3 < \dfrac{M}{3^{2k-1}} < 5.88\]
である。$M$が奇数であることを考慮して$M = 5\cdot 3^{2k-1}$でなければならない。しかし, $M=m^2 - m +1$は決して$5$で割れないことが$\bmod{5}$の考察によって分かるので不適である\footnote{たとえば, $4M = (2m-1)^{2} + 3$を見て$\bmod{5}$で$2$は平方剰余でないことからすれば明らかである。}。\\
\step{3}{TypeII, $p\neq 3$の場合}\\
$N = m^2 + m +1$とおく。このとき不定方程式は
\[p^{3k} + p^{2k}m = (m-1)N\]
である。$N$は常に奇数である。$\mr{gcd}(m-1,N) = 1,3$, $p\neq 3$から同様な理由によって$m-1$, $N$のいずれか一方が$p^{2k}$で割れる。$p^{k} \leq m-1 < 1.4 p^{k} < p^{2k}$より$m-1$が$p^{2k}$で割れることはない。よって$N$は$p^{2k}$で割れる。$p^{k} < m < 1.4p^{k}$より$p^{2k} < N < 3p^{2k}$なので\footnote{$1.4p^{k} + 1 < p^{2k}$より$N < 1,96p^{2k} + 1.4p^{k} + 1<2.96p^{2k}$}$N = 2p^{2k}$でなければならないが, $N$は奇数なので不適。\par
\step{4}{TypeII, $3^{3k} + 3^{2k} m - m^3 = -1$}\\
$3^{2k}(3^{k} + m) = (m-1)N$を考察する。同様の理由で$m-1,N$は両方$3$で割れる。このため. $3^{k} + m$は$3$で割れないので左辺は$3^{2k+1}$では割れない。\par
公約数の制約により$m-1,N$のいずれか一方が"ちょうど"$3^{2k-1}$で割れる。$k\geq 2$であると$3^{k}\leq m-1 < 1.4\cdot  3^{k} < 3^{2k-1}$なので$m-1$は$3^{2k-1}$では割れない。$k=1$の場合, $3<m<4.2$であって, $m=4$である($3$で割れない整数)。このとき,
\[27 + 9\cdot 4 = 63 = 4^{3} - 1\]
なので方程式を満たしている。このときもちろん$f\parena{\dfrac{4}{3}} = -27$で整数となるのでこれらは条件に適している。\par
$N$が$3^{2k-1}$で割れる場合を考える。$3^{k} < m < 1.4 \cdot 3^{k}$より$3^{2k} < N \leq 2.96\cdot 3^{2k}$である\footnote{$m^2 + m + 1 < 1.96\cdot 3^{2k} + 1.4\cdot 3^{k} + 1<1.96\cdot 3^{2k} + 3^{2k}$ より。}。よって
\[3 < \dfrac{N}{3^{2k-1}} < 8.88\]
なので$\dfrac{N}{3^{2k-1}} \in \{ 4,5,6,7,8 \}$である。$N$が奇数であること, このStepの最初に述べた理由から$N$が$3$で割れないこと, および$N\not\equiv 0\pmod{5}$であることから$N=7\cdot 3^{2k-1}$でなければならない。元の式に戻ると
\[3(3^{k} + m) = (m-1)\dfrac{N}{3^{2k-1}} = 7(m-1)\]
より, 整理して$3^{k+1} + 7 = 4m$を得る。よって,
\[112\cdot 3^{2k-1} = 16N = (4m+2)^2 + 12 = (3^{k+1} + 9)^{2} + 12\]
である。最右辺は$9$の倍数ではない。$k\geq 2$であると最左辺が$9$の倍数なので不適である。$k=1$の場合はすでに考察している。\par
以上より,
\[\bolm{(p,m,k) = (3,4,1)}\]
